\ticket{26}{Процессная модель управления проектом: основные процессы и их взаимодействие. Особенности управления программным проектом}

\begin{examplan}
    \item Дать определение управления проектом и процессной модели.
    \item Перечислить группы процессов управления проектом.
    \item Описать взаимодействие процессов.
    \item Назвать особенности управления программным проектом.
\end{examplan}

\subsection*{Короткая версия для письменного ответа}

\textbf{Управление проектом} --- применение знаний, навыков, инструментов и методов к проектной деятельности для достижения целей проекта. \textbf{Процессная модель} рассматривает управление проектом как набор взаимосвязанных процессов с входами, инструментами, методами и выходами.

\subsection*{Группы процессов}

\textbf{Инициация} формально запускает проект или фазу. Результаты: устав проекта, назначение менеджера, определение заинтересованных сторон.

\textbf{Планирование} уточняет цели и способы достижения результата. Включает сбор требований, определение содержания, WBS, расписание, бюджет, качество, ресурсы, коммуникации, риски, закупки и план работы со стейкхолдерами.

\textbf{Исполнение} выполняет работы по плану: создание результатов, управление командой, качеством, коммуникациями, знаниями, закупками и вовлечением заинтересованных сторон.

\textbf{Мониторинг и контроль} отслеживает фактическое исполнение, сравнивает его с планом, выявляет отклонения, управляет изменениями, контролирует содержание, сроки, стоимость, качество, риски и ресурсы.

\textbf{Закрытие} формально завершает проект или фазу: приемка результатов, архивирование, закрытие контрактов, фиксация извлеченных уроков.

\subsection*{Взаимодействие процессов}

Группы процессов не являются строго последовательными фазами. Планирование, исполнение и контроль часто идут параллельно и повторяются итеративно. Контроль порождает запросы на изменения, которые могут возвращать проект к перепланированию.

Типовая логика:
\begin{enumerate}
    \item инициация дает основание проекта;
    \item планирование задает базовые планы;
    \item исполнение создает продукт;
    \item мониторинг сравнивает факт с планом;
    \item изменения проходят контроль и обновляют планы;
    \item закрытие фиксирует результат и опыт.
\end{enumerate}

\subsection*{Особенности программного проекта}

Управление разработкой ПО имеет специфику:
\begin{itemize}
    \item требования часто меняются по мере понимания задачи;
    \item трудоемкость и сроки трудно оценивать точно;
    \item результат нематериален и сложен для визуальной проверки;
    \item качество требует тестирования на разных уровнях;
    \item велико влияние квалификации и коммуникации команды;
    \item важны управление версиями, конфигурациями и зависимостями;
    \item возможны технический долг и архитектурные риски;
    \item часто применяются итеративные и гибкие подходы: Scrum, Kanban, инкрементальная разработка.
\end{itemize}

\textbf{Программный продукт} включает программы, процедуры, правила, документацию и данные, поставляемые пользователю.

\textbf{Процесс разработки ПО} обычно включает требования, проектирование, кодирование, тестирование, развертывание, сопровождение и обновления.

\begin{important}
    \item Процессная модель описывает управление проектом через группы процессов.
    \item Основные группы: инициация, планирование, исполнение, мониторинг и контроль, закрытие.
    \item Процессы взаимодействуют итеративно, особенно через управление изменениями.
    \item В программных проектах особенно важны требования, качество, команда, версии и изменения.
    \item Agile-подходы полезны при высокой неопределенности и изменчивости требований.
\end{important}
